\documentclass[fontsize=10pt]{scrreprt}
% \usepackage[a4paper, left=2cm, right=2cm, top=2cm, bottom=2cm]{geometry}

% \usepackage{import}          
% \usepackage{amsmath, amssymb, amsthm}
% \usepackage{commath}
% \usepackage{graphicx}
% \usepackage{multicol}        
% \usepackage{tcolorbox}       
% \usepackage{hyperref}        
% \usepackage{fancyhdr}
% \usepackage[acronym]{glossaries}

% \usepackage{booktabs}
% \usepackage{colortbl}     

% \usepackage{fancyhdr}

% \usepackage{pgffor}

% \usepackage{amsthm}

% \tcbuselibrary{theorems}

% Eigene Stylings
\usepackage{styling/settings}
\usepackage{styling/style}
\usepackage{styling/commands/math_short_commands}
\usepackage{styling/commands/general_short_commands}
\usepackage{styling/commands/default-expressions}
\usepackage{styling/info-boxen}

\title{Ana II Transkiption}
\author{Finn Zeumer }
\date{Mai 2025}

\begin{document}

    \begin{titlepage}
    \begin{center}
        \vspace*{3cm}
        
        \Huge\textbf{Analisys II Transkiption} \\
        \large{Differential- und Integralrechnung in mehreren Variablen}
        
        \vspace{1.5cm}

        \large{Universität Heidelberg} \\
        \small{Dozent: Ekaterina Kostina }
        
        \vspace{7cm}
        
        \large{Finn Zeumer \\
        Matrikelnummer: 4749114}
        
        \vfill
        
        Stand: {\today}
        
    \end{center}
\end{titlepage} 

    % Römische Zahlen für Content bis zum Inhaltsverzeichnis
    \pagenumbering{roman}
    \chapter{Vorwort}
% \addcontentsline{toc}{chapter}{A. Vorwort}
% \chaptermark{Vorwort}

\section*{Generelles}
\addcontentsline{toc}{section}{A.1 Generelles}
Dieses Skript entstand im Rahmen meiner Vorlesung zur Analysis 2, die ich als Physikstudent im Sommersemester 2026 besuche/besucht habe. Ziel dieses Skripts ist es, die wesentlichen Inhalte der Vorlesung systematisch und verständlich zusammenzufassen, um sowohl die theoretischen Grundlagen als auch die praktische Anwendung der behandelten Konzepte greifbar zu machen.

Als Physikstudent verfolge ich bei der Aufarbeitung der Inhalte eine andere Motivation als es ein reiner Mathematiker tun würde. Mein Schwerpunkt liegt weniger auf formaler Vollständigkeit und mathematischer Rigorosität, sondern vielmehr auf einem anschaulichen Verständnis und der Anwendbarkeit der Methoden im physikalischen Kontext (und dass ich die Klausuren bestehe). Aus diesem Grund sind zahlreiche Beispiele, Skizzen und Visualisierungen enthalten, die das Verständnis der abstrakten Konzepte erleichtern sollen und häufig einen eindeutigen Bezug zur Physik haben.

Alle Grafiken und Skizzen in diesem Skript habe -- wenn nicht anders vermerkt -- ich selbst erstellt (häufig auch nachbildungen aus den Vorlesungsmaterials). Ich habe mich bemüht, die Inhalte sorgfältig und nachvollziehbar zusammenzustellen. Dennoch kann ich weder Vollständigkeit noch absolute Richtigkeit garantieren. Dieses Skript stellt keine offizielle Vorlesungsmitschrift dar, sondern spiegelt meine persönliche Auseinandersetzung mit dem Stoff wider. Fehler und Ungenauigkeiten sind daher nicht ausgeschlossen.

Ich hoffe, dass dieses Skript als hilfreiche Ergänzung zur Vorlesung dienen kann und sowohl mir selbst als auch Kommiliton*innen beim Lernen und Wiederholen der Inhalte von Nutzen ist.

\section*{Vorlesungsinformationen}
Die Vorlesung im Sommersemester wird von \emph{Frau Prof. Dr. Ekaterina A. Kostina} gehalten. Die Folien können auf \href{https://mampf.mathi.uni-heidelberg.de/lectures/313}{MaMpf} abgerufen werden.
Da ich von Recht keine Ahnung habe, bitte ich darum, dieses Material nicht weiter zu teilen, da ich nicht weiß, in wie weit der Inhalt geistiges Eigentum der Dozentin ist. Ich möchte betonen, dass ich explizit kein Anspruch auf geistiges Eigentum bei diesem Werk erhebe.

\section*{Notation}
\addcontentsline{toc}{section}{A.2 Notation}
Es wird sich an der Notation der Folien orientiert, jedoch werden einige Formulierungen hier besonders betont, die das Verständnis stärken soll. 
Zudem wird immer wieder

\subsection*{Kompleze Einheit}
Die Komplexe Einheit $i$ soll hier in diesem Transskipt blau hervorgehoben werden, um die Klarheit für komplexwertige Ausdrücke zu steigern. Es wird also notiert:
$$
    \i
$$

\subsection*{Normen}
In dem Skript werden häufig Normen genannt und genutzt. Diese werden normaler Weise mit 
$$
    \| \cdot \|_\text{Norm}
$$
benannt. Für die $L^p$-Normen gilt somit:
$$
    \| \cdot \|_{L^p} \equiv \| \cdot \|_p 
$$
Zudem wird der sonderfall der $L^2$-Norm betrachtet:
$$
    \| \cdot \|_{L^2} \equiv \| \cdot \|_2  \equiv \| \cdot \| 
$$

    \renewcommand{\contentsname}{Inhaltsverzeichnis}
    \tableofcontents

    \newpage

    % Seitenzahl auf 1 Setzen für textlichen Inhalt
    \pagenumbering{arabic}

    % -------------------
    % Kapitel 1: Topologische Grundlagen
    % -------------------
    \section{Topologische Grundlagen}

    \foreach \file in {
        Einleitung,
        Metrische_Raeume,
        Konvergenz_und_Stetigkeit,
        Kompaktheit,
        Zusammenhang
    }{
        \input{pages/chapter_25/Kapitel_1-Topologische_Grundlagen/\file}
    }

    \newpage

    % -------------------
    % Kapitel 2: Kurven
    % -------------------
    \section{Kurven in $\mathbb{R}^n$}

    \foreach \file in {
        Einleitung,
        Regularitaet_und_Tangentialvektor,
        Rektifizierbarkeit_und_Bogenlaenge
    }{
        \input{pages/chapter_25/Kapitel_2-Kurven/\file}
    }

    \newpage

    % -------------------
    % Kapitel 3: Differentialrechnung in mehreren Variablen
    % -------------------
    \section{Differentialrechnung in mehreren Variablen}

    \foreach \file in {
        Partielle_Ableitungen_und_Richtungsableitungen,
        Differenzierbarkeit,
        Hoehere_Ableitungen,
        Vertauschen_von_Differentiation_und_Integration,
        Lokale_Extrema
    }{
        \input{pages/chapter_25/Kapitel_3-Differentialrechnung_in_mehreren_Variablen/\file}
    }

    \newpage

    % -------------------
    % Kapitel 4: Gleichungen und Mannigfaltigkeiten 
    % -------------------
    \section{Gleichungen und Mannigfaltigkeiten}

    \foreach \file in {
        Einleitung_K4,
        Der_Banachsche_Fixpunktsatz,
        Lokale_Diffeomorphismen,
        Implizite_Funktionen_und_Untermannigfaltigkeiten
    }{
        \input{pages/chapter_25/Kapitel_4-Gleichungen_und_Mannigfaltigkeiten/\file}
    }
\end{document}
